\section{The corrected rank of the regional frustration-free truncation}

\begin{lemma}[Support of complementary products]
    \label{thm:physical_complementary_product_support}
    \lean{QuantumCircuit.mem_supportedOperators_complementProd}
    \leanok
    Let $A_0,\ldots,A_{n-1}$ be operators supported on a common set
    of sites $S$. Then $(I-A_{n-1})\cdots(I-A_0)$ is supported on $S$.
\end{lemma}
\begin{proof}
    \leanok
    The identity belongs to the supported-operator algebra, which is
    closed under subtraction and multiplication. Induct on $n$.
\end{proof}

\begin{lemma}[Support of the actual interval parity products]
    \label{thm:physical_regional_parity_support}
    \lean{TNLean.MPS.AreaLaw.isSupportedOn_physicalParityProduct,
        TNLean.MPS.AreaLaw.isSupportedOn_regionalParityProduct}
    \leanok
    \uses{def:physical_ff_regional_truncation,thm:physical_complementary_product_support}
    Restriction to the bonds internal to $J=[a,b)$ preserves each
    designated bond support. The corresponding actual parity products
    $Q_0,Q_1$ are supported on $J$. More generally, a parity product is
    supported on any common support of its selected terms.
\end{lemma}
\begin{proof}
    \leanok
    Every retained bond $\{i,i+1\}$ is contained in $J$; omitted terms
    vanish. Apply the preceding product-support theorem. No projection,
    commutation or positivity assumption is needed.
\end{proof}

\begin{lemma}[Ranks of the two actual parity products]
    \label{thm:physical_parity_product_cut_rank}
    \lean{TNLean.MPS.AreaLaw.operatorSchmidtRank_physicalParityProduct_le_sq,
        TNLean.MPS.AreaLaw.operatorSchmidtRank_physicalParityProduct_le_one,
        TNLean.MPS.AreaLaw.sum_operatorSchmidtRank_physicalParityProduct_le}
    \leanok
    \uses{thm:physical_regional_parity_support}
    Let $q\ge1$ and let $h_i$ be arbitrary nearest-neighbour supported
    operators. At any site cut, each actual parity product has operator
    Schmidt rank at most $q^2$. The parity not containing the possible
    crossing bond has rank at most one. Consequently
    \begin{align*}
        \operatorname{OSR}(Q_0)+\operatorname{OSR}(Q_1)\le q^2+1.
    \end{align*}
\end{lemma}
\begin{proof}
    \leanok
    Follow the actual ordered product up to the cut. Every noncrossing
    complementary factor has rank at most one. Before the unique possible
    crossing bond, their product therefore has rank at most one. The
    crossing complementary factor has rank at most $q^2$, by its two-site
    support. Every later factor again has rank at most one. The other
    parity contains no crossing factor. Apply the product inequality
    at each step. Cuts at the endpoints contain no crossing bond.
\end{proof}

\begin{lemma}[Corrected cut rank of the actual regional truncation]
    \label{thm:physical_ff_regional_corrected_rank}
    \lean{TNLean.MPS.AreaLaw.operatorSchmidtRank_physicalFrustrationFreeTruncation_le}
    \leanok
    \uses{def:physical_ff_regional_truncation,thm:physical_regional_parity_support,
        thm:physical_parity_product_cut_rank}
    For arbitrary nearest-neighbour supported operators, the actual
    regional truncated operator satisfies, at every cut,
    \begin{align*}
        \operatorname{OSR}(\widetilde H_J)\le q^2+3.
    \end{align*}
    This includes empty intervals, endpoint cuts, zero local dimension
    and the empty chain. It corrects the false additive constant in
    \cite[Section~6.1]{Arad2016RigorousRG}; the printed $q^2+2$ assertion
    and its nondegenerate counterexample are recorded in
    \cite{gap:alvv16_frustration_free_truncation_rank}.
\end{lemma}
\begin{proof}
    \leanok
    Split the actual exterior bond sum into its left, right and possible
    crossing terms. If no exterior term crosses the cut, combine $2I$
    with the left sum. These two one-sided terms have total rank at most
    two. The two actual parity products contribute at most $q^2+1$.
    Hence their sum has rank at most $q^2+3$.

    If an exterior bond crosses, $J$ lies entirely on one side of that
    cut. Both actual parity complements can be combined with that side's
    exterior sum. The two one-sided sums contribute at most two and the
    crossing bond at most $q^2$, giving the stronger $q^2+2$ bound in this
    case. No decomposition or rank estimate is assumed: each piece is
    the corresponding actual bond sum or parity product.

    For $q=0$ on a nonempty chain the configuration space is empty and
    every matrix vanishes. On the empty chain the original sum vanishes
    and both parity products are the identity, so $\widetilde H_J=0$.
\end{proof}
